% ============================================================================
% BLOCK-003 Derivation v9: NC-2 Attempt (DGP / Induced Gravity)
% ============================================================================
\documentclass[11pt,a4paper]{article}

\usepackage{fontspec}
\usepackage{amsmath,amssymb,amsthm}
\usepackage{physics}
\usepackage{booktabs}
\usepackage{array}
\usepackage{xcolor}
\usepackage[hidelinks]{hyperref}
\usepackage[margin=2.5cm]{geometry}
\usepackage{fancyhdr}
\usepackage{tcolorbox}

% Theorem environments
\newtheorem{proposition}{Proposition}
\theoremstyle{definition}
\newtheorem{definition}{Definition}

% Epistemic tags
\newcommand{\BL}{\textsc{[BL]}}
\newcommand{\Post}{\textsc{[P]}}
\newcommand{\Deriv}{\textsc{[Dc]}}
\newcommand{\Open}{\textsc{[Open]}}

% Boxes
\newtcolorbox{outcomebox}[1][]{colback=green!5,colframe=green!50!black,
  title=Outcome,fonttitle=\bfseries,#1}
\newtcolorbox{statusbox}[1][]{colback=yellow!10,colframe=orange!80!black,
  title=#1,fonttitle=\bfseries}

% Header/footer
\pagestyle{fancy}
\fancyhf{}
\fancyhead[L]{\small BLOCK-003 / Derivation v9}
\fancyhead[R]{\small EDC Gravity Program}
\fancyfoot[C]{\thepage}

\title{\textbf{EDC BLOCK-003 (Derivation v9):}\\[0.3em]
NC-2 Attempt: DGP / Induced Gravity\\[0.3em]
\large Can Brane-Localized Curvature Break the Degeneracy?}
\author{EDC Working Document\\[0.5em]
\small 2026-02-02 --- Pivot attempt after NC-1 inconclusive}
\date{}

\begin{document}
\maketitle

\begin{abstract}
Derivation v8 showed that NC-1 (zero-mode normalization) yields an identity
rather than an independent equation: $C$ and $\sigma$ always appear in
the same combination. This note tests NC-2: Dvali-Gabadadze-Porrati (DGP)
induced gravity, where a brane-localized Einstein-Hilbert term provides
an independent 4D Planck mass $M_4$ distinct from the bulk-derived contribution.
\textbf{Outcome:} INCONCLUSIVE. DGP introduces a new scale $r_c = M_5^3/M_4^2$
(crossover radius), but $M_4^2$ itself depends on unknown parameters
($N$, $\Lambda_{\rm UV}$, regularization). The degeneracy is not broken;
it is relocated from $C$ to the induced term coefficient.
BLOCK-003 remains \Open.
\end{abstract}

\tableofcontents
\newpage

% ============================================================================
\section{Motivation: Why NC-2 After NC-1?}
% ============================================================================

From v8, NC-1 (canonical KK reduction) produced:
\begin{equation}
G_N = \frac{C \sigma^{-3/4}}{8\pi R_\xi}
\end{equation}
where $C$ and $\sigma$ are degenerate --- they always appear together.
This is a \textit{reparametrization}, not an independent equation.

\textbf{NC-2 hypothesis:} In DGP-type models, 4D gravity emerges from a
brane-localized curvature term, providing an \textit{independent} scale
$M_4$ that could break the $C$--$\sigma$ degeneracy.

% ============================================================================
\section{DGP Setup}
% ============================================================================

\subsection{Action}

The DGP action combines bulk and brane Einstein-Hilbert terms:
\begin{equation}
S = \underbrace{\frac{M_5^3}{2} \int_{\mathcal{M}^5} d^5X \sqrt{-g^{(5)}} R^{(5)}}_{S_{\rm bulk}}
+ \underbrace{\frac{M_4^2}{2} \int_{\Sigma^4} d^4x \sqrt{-h} R^{(4)}}_{S_{\rm brane}}
+ S_{\rm matter}
\label{eq:dgp_action}
\end{equation}
where:
\begin{itemize}
\item $M_5^3 = 1/\kappa_5^2$ is the 5D Planck mass cubed
\item $M_4^2$ is the 4D brane-localized Planck mass squared
\item $R^{(5)}$, $R^{(4)}$ are 5D and 4D Ricci scalars
\end{itemize}

\subsection{Key Feature: Two Independent Scales}

Unlike pure KK reduction (NC-1), DGP has \textit{two} independent mass scales:
\begin{itemize}
\item $M_5$: bulk gravitational scale
\item $M_4$: brane gravitational scale
\end{itemize}

The crossover radius is:
\begin{equation}
r_c = \frac{M_5^3}{M_4^2}
\label{eq:rc}
\end{equation}

At distances $r \ll r_c$: 4D gravity dominates ($G_N \sim 1/M_4^2$).

At distances $r \gg r_c$: 5D gravity dominates (leakage into bulk).

% ============================================================================
\section{Effective 4D Newton Constant}
% ============================================================================

\subsection{DGP Graviton Propagator}

The effective 4D gravitational coupling at distance $r$ is:
\begin{equation}
G_N^{\rm eff}(r) = \frac{1}{8\pi M_4^2} \cdot F(r/r_c)
\end{equation}
where $F(x) \to 1$ as $x \to 0$ (4D regime).

In the 4D regime ($r \ll r_c$):
\begin{equation}
\boxed{G_N = \frac{1}{8\pi M_4^2}}
\label{eq:gn_dgp}
\end{equation}

\subsection{Comparison to NC-1}

In NC-1 (pure KK):
\begin{equation}
G_N^{\rm NC1} = \frac{\kappa_5^2}{8\pi L} = \frac{1}{8\pi M_5^3 L}
\end{equation}

In NC-2 (DGP):
\begin{equation}
G_N^{\rm NC2} = \frac{1}{8\pi M_4^2}
\end{equation}

\textbf{Key difference:} In DGP, $G_N$ is determined by $M_4$, not by $M_5$
and $L$. This \textit{could} break the degeneracy if $M_4$ is independently
specified.

% ============================================================================
\section{Where Does $M_4$ Come From?}
% ============================================================================

\subsection{Induced Gravity Mechanism}

In Sakharov/Adler induced gravity, $M_4^2$ arises from quantum loops of
brane-localized matter:
\begin{equation}
M_4^2 = \frac{N}{16\pi^2} \Lambda_{\rm UV}^2 \cdot f(\text{reg})
\label{eq:m4_induced}
\end{equation}
where:
\begin{itemize}
\item $N$: number of matter species on the brane
\item $\Lambda_{\rm UV}$: UV cutoff
\item $f(\text{reg})$: regularization-dependent $\mathcal{O}(1)$ factor
\end{itemize}

\subsection{EDC Interpretation}

For EDC, possible identifications:
\begin{itemize}
\item $\Lambda_{\rm UV} = \sigma^{1/4}$ (tension-derived cutoff) \Post
\item $\Lambda_{\rm UV} = R_\xi^{-1}$ (weak scale cutoff) \Post
\item $N = ?$ (unknown; EDC does not specify species count)
\end{itemize}

With $\Lambda_{\rm UV} = \sigma^{1/4}$:
\begin{equation}
M_4^2 \sim \frac{N}{16\pi^2} \sigma^{1/2}
\end{equation}

Then:
\begin{equation}
G_N = \frac{1}{8\pi M_4^2} \sim \frac{16\pi^2}{8\pi N \sigma^{1/2}}
= \frac{2\pi}{N \sigma^{1/2}}
\label{eq:gn_dgp_sigma}
\end{equation}

% ============================================================================
\section{Does NC-2 Break the Degeneracy?}
% ============================================================================

\begin{proposition}[NC-2 Degeneracy Relocation]
In DGP/induced gravity, the unknown constant $C$ from NC-1 is replaced by
the unknown species count $N$ (and regularization factor $f$). The degeneracy
is relocated, not eliminated.
\end{proposition}

\begin{proof}
From Eq.~\eqref{eq:gn_dgp_sigma}:
\[
G_N \sim \frac{2\pi}{N \sigma^{1/2}}
\]
To predict $G_N$, we need $N$ and $\sigma$. Neither is specified by EDC:
\begin{itemize}
\item $\sigma$: brane tension (same unknown as NC-1)
\item $N$: species count (new unknown, replaces $C$)
\end{itemize}
The product $N \sigma^{1/2}$ is degenerate, just as $C \sigma^{3/4}$ was
in NC-1.
\end{proof}

\subsection{Comparison of Degeneracies}

\begin{table}[h]
\centering
\begin{tabular}{lll}
\toprule
\textbf{Approach} & \textbf{$G_N$ Formula} & \textbf{Unknown Combination} \\
\midrule
NC-1 (KK) & $C \sigma^{-3/4} / (8\pi R_\xi)$ & $C \sigma^{-3/4}$ \\
NC-2 (DGP) & $2\pi / (N \sigma^{1/2})$ & $N \sigma^{1/2}$ \\
\bottomrule
\end{tabular}
\caption{Both approaches have one degenerate combination.}
\end{table}

% ============================================================================
\section{What Would Fix NC-2?}
% ============================================================================

For NC-2 to close BLOCK-003, EDC would need to specify:

\begin{enumerate}
\item \textbf{Species count $N$:} How many field species live on the brane?
EDC's Z$_6$ structure suggests discrete particle content, but $N$ is not
derived.

\item \textbf{Regularization scheme:} The factor $f(\text{reg})$ is
scheme-dependent. A canonical choice (e.g., dimensional regularization)
would fix this to $\mathcal{O}(1)$.

\item \textbf{Brane tension $\sigma$:} Same requirement as NC-1.
\end{enumerate}

\textbf{Minimal fix:} If $N$ is small and known (e.g., $N = 1$ for a
single Plenum field), and $\sigma$ is specified, then $G_N$ becomes
predictable.

% ============================================================================
\section{Outcome}
% ============================================================================

\begin{outcomebox}
\textbf{Status: INCONCLUSIVE}

\textbf{What was achieved:}
\begin{itemize}
\item DGP framework applied to EDC \Deriv
\item $G_N = 2\pi / (N \sigma^{1/2})$ derived (with $\Lambda_{\rm UV} = \sigma^{1/4}$)
\item Identified that degeneracy is \textit{relocated}, not eliminated
\end{itemize}

\textbf{What was not achieved:}
\begin{itemize}
\item $N$ is not specified by EDC
\item Regularization factor remains $\mathcal{O}(1)$ uncertainty
\item $\sigma$ still required (same as NC-1)
\end{itemize}

\textbf{Conclusion:}
NC-2 does not break the degeneracy. The unknown $C$ is replaced by unknown
$N$. BLOCK-003 remains \Open.
\end{outcomebox}

% ============================================================================
\section{Implications for BLOCK-003}
% ============================================================================

\begin{statusbox}[Diagnosis]
Both NC-1 and NC-2 fail for the same structural reason: EDC lacks a
\textbf{normalization anchor}---a first-principles way to fix one
dimensionless number ($C$, $N$, or equivalent).

\textbf{Possible interpretations:}
\begin{enumerate}
\item \textbf{Genuine underdetermination:} The 5D framework inherently
has one free parameter that must be fixed by observation (calibration).
\item \textbf{Missing physics:} A stabilization mechanism (modulus fixing,
topological quantization) is needed but not present in current EDC.
\item \textbf{Tautology:} The derivation is circular---4D gravity was
presupposed somewhere in the axioms.
\end{enumerate}

\textbf{Recommendation:} Perform tautology audit (derivation v10) before
attempting further normalization candidates.
\end{statusbox}

\end{document}
